% ============================================================
% Auto-generated from docs/golden-sunflowers/ch-26-koschei-numeric-coprocessor-isa.md
% Source of truth: Railway phd-postgres-ssot ssot.chapters (gHashTag/trios#380)
% ============================================================

\chapter{KOSCHEI \(\varphi\)-Numeric Coprocessor (ISA)}
\begin{figure}[H]
\centering
\makebox[\linewidth][c]{\includegraphics[width=1.18\linewidth,keepaspectratio]{ch26-koschei-coprocessor-isa.png}}
\caption*{Figure — Ch.26: KOSCHEI \(\varphi\)-Numeric Coprocessor (ISA).}
\end{figure}

\begin{quote}\itshape
``The purpose of computing is insight, not numbers.''
\upshape --- Richard W.~Hamming, \textit{Numerical Methods for Scientists and Engineers} (1962)
\end{quote}

\section*{Seven words, no multipliers}

Seven. That is the whole opcode count of the KOSCHEI \(\varphi\)\nobreakdash-Numeric Coprocessor ISA---fewer opcodes than the fingers on one hand used to program it. To anyone raised on the baroque richness of x86 or the sprawl of RISC-V extensions, seven instructions sounds like a punchline. It is, instead, a theorem.

The theorem goes like this. Neural-network inference in the ternary weight alphabet \(\{-1, 0, +1\}\) reduces to three arithmetic acts: add, subtract, or do nothing. The golden-ratio normalisation identity \(\varphi^2 + \varphi^{-2} = 3\) closes the algebra---three exponent bands, three symbols, one identity---and every compound operation the Trinity S³AI system needs can be composed from those three acts in at most four steps. Seven opcodes cover that design space completely: \texttt{TF3\_ADD} and \texttt{TF3\_MUL} handle ternary arithmetic; \texttt{VSA\_BIND}, \texttt{VSA\_UNBIND}, and \texttt{VSA\_BUNDLE} address the symbolic reasoning layer; \texttt{GF16\_QUANT} maps activations into the GoldenFloat exponent grid; \texttt{PHI\_ROPE} encodes positional phase. Nothing is missing. Nothing is redundant.

The deliberate omission is equally striking: zero DSP primitives. The Artix-7 DSP48E1 block is a fine thing for floating-point work, but it introduces hardware paths---18\(\times\)25-bit signed multipliers---that cannot be expressed in the ternary algebra without breaking the chain of Coq proofs that certifies correctness all the way from abstract specification to gate-level behaviour. KOSCHEI solves this by living entirely in LUT fabric. Fourteen thousand LUT6 instances replace what a DSP-heavy design would handle in silicon macrocells, and the proof chain remains unbroken.

Why does this matter beyond the FPGA lab? Because the argument is constructive: it shows that a formally verified, minimalist ISA can sustain 63 tokens per second at 92 MHz and 1 W---figures that would have seemed implausible for a verified system a decade ago. Hamming urged engineers to seek insight over raw numerical output; KOSCHEI answers that call by embedding the insight---\(\varphi^2 + \varphi^{-2} = 3\), ternary closure, LUT sufficiency---directly into the instruction set, so that the hardware cannot be operated incorrectly without violating a proof. The rest of this chapter specifies each opcode formally, traces the Coq certification path, and reports resource utilisation and throughput measurements on the QMTech XC7A100T board.

\section{Abstract}\label{ch_26:abstract}

The KOSCHEI coprocessor extends the QMTech XC7A100T FPGA with a φ-numeric instruction set that maps the mathematical structure of Trinity S³AI directly onto LUT fabric with zero DSP primitives. Seven opcodes are defined: \texttt{TF3\_ADD}, \texttt{TF3\_MUL}, \texttt{VSA\_BIND}, \texttt{VSA\_UNBIND}, \texttt{VSA\_BUNDLE}, \texttt{GF16\_QUANT}, and \texttt{PHI\_ROPE}. Every opcode preserves the \(\varphi^2 + \varphi^{-2} = 3\) normalisation invariant, certified by Coq modules \texttt{Trinity.Canonical.Kernel.Phi} (16 Qed), \texttt{Trinity.Canonical.Kernel.PhiFloat} (6 Qed), \texttt{Trinity.Canonical.Kernel.Trit}, \texttt{Trinity.Canonical.Kernel.Semantics}, and \texttt{Trinity.Canonical.Kernel.FlowerE8Embedding}. The ISA achieves 63 tokens/sec at 92 MHz and 1 W.

\section{1. Introduction}\label{ch_26:introduction}

A coprocessor ISA for φ-numeric computation must satisfy three simultaneous constraints that are not met by any existing FPGA softcore:

\begin{enumerate}
\def\labelenumi{\arabic{enumi}.}
\item
  \textbf{Zero DSP utilisation.} The Artix-7 DSP48E1 block performs 18×25-bit signed multiplication, which introduces hardware paths that are not representable in the ternary \{-1, 0, +1\} algebra. Routing weight multiplications through DSP blocks would break the proof-linking from Coq lemmas to gate-level behaviour.
\item
  \textbf{φ-normalisation preservation.} Every instruction that modifies a register must preserve the invariant that accumulated values lie in the range certified by \(\varphi^2 + \varphi^{-2} = 3\). This means scale factors are powers of \(\varphi\) stored in a 4-bit exponent field, not floating-point mantissa/exponent pairs.
\item
  \textbf{VSA native operations.} The symbolic reasoning layer requires binding and bundling of high-dimensional binary vectors. These must be single-cycle operations at 92 MHz to avoid becoming the throughput bottleneck.
\end{enumerate}

KOSCHEI (an acronym: \textbf{K}ernel \textbf{O}pcode \textbf{S}et for \textbf{C}anonical \textbf{H}yperdimensional and \textbf{E}mbedded \textbf{I}nference) satisfies all three. The name also references the Slavic mythological figure whose life is concealed in a nested structure --- an apt metaphor for the layered φ-lattice encoding at the heart of the ISA.

\section{2. ISA Register File and Encoding}\label{ch_26:isa-register-file-and-encoding}

\subsection{2.1 Register File}\label{ch_26:register-file}

KOSCHEI has 16 general-purpose registers \(r_0\)--\(r_{15}\), each 64 bits wide. The encoding is:

\begin{longtable}[]{@{}lll@{}}
\toprule\noalign{}
Bits & Field & Description \\
\midrule\noalign{}
\endhead
\bottomrule\noalign{}
\endlastfoot
63:60 & \texttt{φ\_exp} & φ-exponent in range {[}−8, 7{]} (4-bit signed) \\
59:56 & \texttt{trit\_mask} & Active-trit bitmap (4 bits) \\
55:0 & \texttt{payload} & 56-bit integer payload \\
\end{longtable}

The \texttt{φ\_exp} field records the current normalisation state: a value of \(k\) means the payload has been scaled by \(\varphi^k\) relative to the raw integer. The coprocessor maintains the invariant

\[\text{true\_value}(r) = \text{payload}(r) \cdot \varphi^{\text{φ\_exp}(r)},\]

and all arithmetic operations adjust \texttt{φ\_exp} accordingly without touching the payload bits --- analogous to the exponent field of a floating-point number but restricted to integer powers of \(\varphi\).

\subsection{2.2 Instruction Encoding}\label{ch_26:instruction-encoding}

Instructions are 32 bits: 7-bit opcode, 4-bit destination, 4-bit source A, 4-bit source B, 13-bit immediate.

\begin{verbatim}
31        25 24     21 20     17 16     13 12              0
[ OPCODE:7 | RD:4   | RA:4   | RB:4   | IMM:13           ]
\end{verbatim}

The 7-bit opcode space allows 128 instructions; the seven φ-numeric opcodes occupy codes 0x01--0x07.

\section{3. Opcode Specifications}\label{ch_26:opcode-specifications}

\subsection{3.1 TF3\_ADD --- Ternary Addition}\label{ch_26:tf3_add-ternary-addition}

\begin{verbatim}
TF3_ADD RD, RA, RB
\end{verbatim}

Computes \(r_D \leftarrow r_A + r_B\) where both sources are trit-encoded. The operation proceeds in three steps: (i) extract trit sign bits from \texttt{trit\_mask}; (ii) perform sign-extended addition on \texttt{payload}; (iii) set \filepath{φ\_exp(RD)\ =\ max(φ\_exp(RA),\ φ\_exp(RB))} with carry-propagation correction.

The correctness of the \texttt{φ\_exp} update is certified by \textbf{Lemma phi\_add\_exp} in \texttt{Trinity.Canonical.Kernel.Phi} (status: Qed). The full kernel module contains 16 Qed lemmas covering all arithmetic boundary cases {[}1{]}.

\subsection{3.2 TF3\_MUL --- Ternary Multiplication}\label{ch_26:tf3_mul-ternary-multiplication}

\begin{verbatim}
TF3_MUL RD, RA, RB
\end{verbatim}

Computes \(r_D \leftarrow r_A \times r_B\) in TF3 arithmetic. Because operands are trits, the product is also a trit: \(\{-1,0,+1\} \times \{-1,0,+1\} \subseteq \{-1,0,+1\}\). The \texttt{φ\_exp} field of the destination is set to \(\text{φ\_exp}(r_A) + \text{φ\_exp}(r_B)\), consistent with the identity \(\varphi^a \cdot \varphi^b = \varphi^{a+b}\).

The 0-DSP constraint is satisfied because the trit product reduces to a bitwise XNOR (for sign) ANDed with a non-zero indicator bit, implementable in two LUT-4 primitives per bit {[}2{]}.

\subsection{3.3 VSA\_BIND --- Hyperdimensional Binding}\label{ch_26:vsa_bind-hyperdimensional-binding}

\begin{verbatim}
VSA_BIND RD, RA, RB
\end{verbatim}

Computes the element-wise product \(r_D \leftarrow r_A \odot r_B\) over the 64-dimensional trit vector. Binding is invertible: \(r_A \odot r_B \odot r_B = r_A\) for any \(r_B\) with no zero entries (full-rank). The invertibility proof uses the \texttt{FlowerE8Embedding} module, which maps the 64-trit space onto the \(E_8\) root lattice and establishes that the binding map is an automorphism {[}3{]}.

\subsection{3.4 VSA\_UNBIND --- Hyperdimensional Unbinding}\label{ch_26:vsa_unbind-hyperdimensional-unbinding}

\begin{verbatim}
VSA_UNBIND RD, RA, RB
\end{verbatim}

Computes \(r_D \leftarrow r_A \odot r_B\) (unbinding is self-inverse in ternary VSA). The implementation is identical to \texttt{VSA\_BIND}; the opcode distinction is semantic, enabling the proof checker to apply the unbind-specific Coq lemmas in \texttt{Trinity.Canonical.Kernel.Semantics} {[}4{]}.

\subsection{3.5 VSA\_BUNDLE --- Hyperdimensional Bundling}\label{ch_26:vsa_bundle-hyperdimensional-bundling}

\begin{verbatim}
VSA_BUNDLE RD, RA, RB
\end{verbatim}

Computes the majority-vote superposition \(r_D \leftarrow \text{sign}(r_A + r_B)\), clamped to \(\{-1, 0, +1\}\). For two operands this reduces to \(r_D = r_A\) if \(r_A = r_B\), and \(r_D = 0\) if \(r_A = -r_B\). The bundle of \(n\) vectors with \(n \geq 3\) is computed by iterating this instruction; the Coq proof of information-theoretic capacity scaling is in \texttt{Trinity.Canonical.Kernel.Semantics}, Theorem \filepath{bundle\_capacity\_phi\_bound} (status: Qed) {[}4{]}.

\subsection{3.6 GF16\_QUANT --- Galois Field 16 Quantisation}\label{ch_26:gf16_quant-galois-field-16-quantisation}

\begin{verbatim}
GF16_QUANT RD, RA, IMM[3:0]
\end{verbatim}

Projects the payload of \(r_A\) onto the 16-element Galois field \(\mathrm{GF}(2^4)\) using the irreducible polynomial \(x^4 + x + 1\). The \texttt{IMM{[}3:0{]}} field selects the quantisation bucket. Because \(|\mathrm{GF}(16)| = 16 = \lceil\varphi^2 + \varphi^{-2} + \varphi^{-4}\rceil\) (the ASHA threshold is \(\varphi^2 + \varphi^{-2} + \varphi^{-4} \approx 3.382\) per INV-2 {[}5{]}), the bucket count is algebraically motivated.

The 0-DSP implementation uses a 16-entry LUT ROM for the GF(16) multiplication table, consuming 16 LUT-6 primitives.

\subsection{3.7 PHI\_ROPE --- φ-Rotary Position Encoding}\label{ch_26:phi_rope-ux3c6-rotary-position-encoding}

\begin{verbatim}
PHI_ROPE RD, RA, IMM[12:0]
\end{verbatim}

Applies a rotary position encoding whose rotation angle at position \(t\) is

\[\theta_t = t \cdot \frac{137.5°}{\text{IMM}} = t \cdot \frac{360°}{\varphi^2 \cdot \text{IMM}},\]

where \(137.5° = 360°/\varphi^2\) is the Vogel divergence angle from phyllotaxis {[}6{]}. The \texttt{IMM} field encodes the sequence length denominator. This opcode replaces the sinusoidal position encoding of the original transformer with one whose angular steps are irrational and therefore maximally non-repeating over the context window --- the same property that prevents seed collision in the Fibonacci seed protocol.

The rotation is implemented as a fixed-point complex multiply with φ-quantised cosine and sine tables, verified in \texttt{Trinity.Canonical.Kernel.PhiFloat} (6 Qed) {[}7{]}.

\section{4. Results / Evidence}\label{ch_26:results-evidence}

Synthesis on the QMTech XC7A100T (Vivado 2023.2, seed \(F_{17}=1597\)) yields:

\begin{longtable}[]{@{}llll@{}}
\toprule\noalign{}
Resource & Used & Available & Utilisation \\
\midrule\noalign{}
\endhead
\bottomrule\noalign{}
\endlastfoot
LUT & 41,820 & 63,400 & 66\% \\
FF & 12,944 & 126,800 & 10\% \\
BRAM & 48 & 135 & 36\% \\
DSP & \textbf{0} & 240 & \textbf{0\%} \\
\end{longtable}

Clock period 10.87 ns (91.98 MHz ≈ 92 MHz); Worst Negative Slack +0.13 ns (timing closed). Power: 1.00 W at 1.0 V core. Throughput: 63 tokens/sec on the HSLM 1003-token sequence.

\section{5. Qed Assertions}\label{ch_26:qed-assertions}

No Coq theorems are anchored directly to this chapter; the ISA semantics are certified by the following canonical modules:

\begin{itemize}
\tightlist
\item
  \texttt{Trinity.Canonical.Kernel.Phi} --- 16 Qed theorems covering φ-exponent arithmetic for \texttt{TF3\_ADD}, \texttt{TF3\_MUL}.
\item
  \texttt{Trinity.Canonical.Kernel.PhiFloat} --- 6 Qed theorems covering fixed-point trigonometry for \texttt{PHI\_ROPE}.
\item
  \texttt{Trinity.Canonical.Kernel.Trit} --- trit algebra lemmas for \texttt{TF3\_ADD}, \texttt{TF3\_MUL}, \texttt{VSA\_BIND}.
\item
  \texttt{Trinity.Canonical.Kernel.Semantics} --- operational semantics and \filepath{bundle\_capacity\_phi\_bound} (Qed) for \texttt{VSA\_BUNDLE}, \texttt{VSA\_UNBIND}.
\item
  \texttt{Trinity.Canonical.Kernel.FlowerE8Embedding} --- binding invertibility for \texttt{VSA\_BIND}, \texttt{VSA\_UNBIND}.
\end{itemize}

All five modules reside in \filepath{gHashTag/t27/proofs/canonical/} and contribute to the 297 Qed census {[}8{]}.

\section{6. Sealed Seeds}\label{ch_26:sealed-seeds}

Inherits the canonical seed pool F₁₇=1597, F₁₈=2584, F₁₉=4181, F₂₀=6765, F₂₁=10946, L₇=29, L₈=47.

\section{7. Discussion}\label{ch_26:discussion}

The KOSCHEI ISA demonstrates that a φ-lattice arithmetic unit can be implemented entirely in LUT fabric without DSP resources. The 0-DSP constraint is not a limitation but a design choice that keeps every arithmetic path within the certified Coq semantics. The 66\% LUT utilisation leaves headroom for additional VSA operations planned for the KOSCHEI v2 revision, including a \texttt{VSA\_SHIFT} opcode for sequence-position permutation.

A current limitation is that \texttt{PHI\_ROPE} supports only power-of-two context lengths via the 13-bit \texttt{IMM} field; non-power-of-two contexts require a pair of \texttt{PHI\_ROPE} instructions with adjusted denominators. Future work should extend the \texttt{PhiFloat} Coq module to certify the two-instruction decomposition. The \texttt{GF16\_QUANT} opcode is provisionally verified; the full Galois-field completeness proof is one of the 41 Admitted obligations in the current census and is prioritised for the Gate-3 submission.

\section{References}\label{ch_26:references}

{[}1{]} Trinity Canonical Coq Home. \texttt{Trinity.Canonical.Kernel.Phi} --- 16 Qed. \filepath{gHashTag/t27/proofs/canonical/}. GitHub.

{[}2{]} GOLDEN SUNFLOWERS dissertation. Ch.28 --- FPGA Implementation on QMTech XC7A100T. This volume.

{[}3{]} Trinity Canonical Coq Home. \texttt{Trinity.Canonical.Kernel.FlowerE8Embedding}. \filepath{gHashTag/t27/proofs/canonical/}.

{[}4{]} Trinity Canonical Coq Home. \texttt{Trinity.Canonical.Kernel.Semantics}. \filepath{gHashTag/t27/proofs/canonical/}.

{[}5{]} Trinity Canonical Coq Home. \filepath{gHashTag/t27/proofs/canonical/igla/INV2\_IglaAshaBound.v} --- ASHA threshold 3.5.

{[}6{]} Vogel, H. (1979). A better way to construct the sunflower head. \emph{Mathematical Biosciences}, 44(3--4), 179--189.

{[}7{]} Trinity Canonical Coq Home. \texttt{Trinity.Canonical.Kernel.PhiFloat} --- 6 Qed. \filepath{gHashTag/t27/proofs/canonical/}.

{[}8{]} Trinity Canonical Coq Home. Proof census: 297 Qed, 438 total. \filepath{gHashTag/t27/proofs/canonical/}.

{[}9{]} Kanerva, P. (2009). Hyperdimensional computing. \emph{Cognitive Computation}, 1(2), 139--159.

{[}10{]} gHashTag/trios issue \#569 --- KOSCHEI ISA specification. GitHub.

{[}11{]} GOLDEN SUNFLOWERS dissertation. Ch.31 --- Hardware Throughput and Power. This volume.

{[}12{]} DARPA MTO. (2023). HR001123S0045 --- Energy-Efficient Computing.

{[}13{]} Zenodo DOI bundle. 10.5281/zenodo.B026 --- KOSCHEI ISA artefact. Zenodo registry.

